\documentclass[11pt]{article}

% --- Required Packages ---
\usepackage{amsmath}
\usepackage{amssymb}
\usepackage{booktabs}
\usepackage{siunitx}
\usepackage{microtype}
\usepackage{float} % Prevents tables from floating
\usepackage{url}

% --- Custom Unit Definitions for siunitx v3 ---
\DeclareSIUnit\pps{pps}

\title{\textbf{Resilient RAP: Autonomous Schema-Drift Remediation Using Hybrid Classical-Quantum Variational Routing Engines}}
\author{\textbf{Tarek Clarke} \\ Tallinn University of Technology (TalTech)}
\date{\monthname\ \the\year}

\begin{document}

\maketitle

\begin{abstract}
Modern Reproducible Analytical Pipelines (RAP) suffer from severe processing bottlenecks when encountering structural schema drift in real-time, high-velocity data streams. Traditional reconciliation requires synchronous CPU branch checking or manual conditional interventions, inducing significant latency penalties. This paper presents an improved iteration of the Resilient RAP Framework, an autonomous orchestration layer that introduces a 24-qubit Variational Quantum Classifier (VQC) routing network to evaluate drift characteristics on-the-fly. By implementing a dual-run pre-warmup phase on AMD Instinct accelerators, we completely isolate HIP kernel compilation cold starts. Our framework demonstrates a reduction in processing latency down to a deterministic $0.437\pm0.005\text{ ms}$ under a $100\%$ schema chaos load, maintaining a $93.58\%$ system accuracy baseline while reducing active GPU resource utilization from $100\%$ to $15\%$.
\end{abstract}

\section{Introduction \& Hardware Environment}
This experiment was tested on the LUMI-G supercomputer, using a single AMD MI250X GPU for each benchmark. While scaling to a multi-GPU benchmark is possible on this platform, single GPU benchmarks were used for sustainability, portability and reproducibility. For comparison, up to 8 GPUs were utilized to assess the utility of an increase in GPU computing power. In order to conduct a quantum routing proof of concept, IBM Cloud's Quantum Open Plan was used, leveraging a 156-qubit IBM Heron r2 QPU using a heavy hex lattice for routing. Ten hours of QPU processing time was secured via EuroHPC on the VQC quantum computer, using a 24-qubit star QPU for extensive benchmarking.

Notably, while routing latency increased from 0~ms to 0.2~ms when switching from GPU-only to Hybrid Quantum Routing, the average reconciliation latency dropped from $\sim$1,250~ms to $\sim$185~ms per 1,000 packets. The throughput experienced a significant increase of 71.5 packets per second (pps) when using Hybrid Quantum Routing, while decreasing the GPU utilization time from 100\% to 15\%. Overall, there was a slight decrease of accuracy from 98.0\% to 94.6\%, but GPU resource usage was significantly decreased, leading to a more sustainable reconciliation process. Since Hybrid Quantum Routing requires training over live compute time during the telemetry ingestion and reconciliation, it is a viable use of resources.

\section{Mathematical Formulations for Resource Profiling}
To establish a rigorous resource evaluation framework for self-healing analytical pipelines, we define the quantitative models governing energy, power, and resilience below.

The cumulative GPU energy consumption $E_{\text{total}}$ is obtained by numerically integrating the time-varying power draw $P(t)$ over the execution duration:
\begin{equation}
E_{\text{total}} = \int_{t_0}^{t_N} P(t) \, dt \approx \sum_{i=0}^{N-1} \frac{P(t_i) + P(t_{i+1})}{2} \Delta t_i
\end{equation}
where $P(t_i)$ is the active power draw sampled at discrete time steps $t_i$, and $\Delta t_i = t_{i+1} - t_i$ is the sampling interval (nominally set to $\qty{0.05}{\second}$).

The normalized energy consumption per $1,000$ anomaly events, denoted as $E_{\text{norm}}$, is computed as:
\begin{equation}
E_{\text{norm}} = \left( \frac{E_{\text{total}}}{N_{\text{drift}}} \right) \times 1000
\end{equation}
where $N_{\text{drift}}$ is the total number of drifted packets that triggered a schema reconciliation stage during the run.

The average active power draw $P_{\text{avg}}$ is calculated as the ratio of total energy to the total execution duration:
\begin{equation}
P_{\text{avg}} = \frac{E_{\text{total}}}{t_N - t_0}
\end{equation}

To quantify the robustness of the self-healing RAP framework under structural schema drift, we utilize the Recovery Coverage metric ($C_{\text{rec}}$), also referred to as the Hosseini Resilience Index:
\begin{equation}
C_{\text{rec}} = \frac{\sum_{i=1}^{n} R_i}{\sum_{i=1}^{n} D_i}
\end{equation}
where $D_i$ represents the number of detected anomalies (drifted parameters) in batch $i$, and $R_i$ is the number of successfully resolved and mapped anomalies in batch $i$.

\section{Experimental Results}

In Table~\ref{tab:chaos_sweep} it is observed that accuracy remains stable, with a lower error range, and latency improves, and energy expenditure increases as drifted packets increase.

\begin{table}[H]
\centering
\small
\caption{BERT Chaos Sweep Results (Mean $\pm$ Std)}
\label{tab:chaos_sweep}
\begin{tabular}{cccccc}
\toprule
Rate (\%) & Total Pkts & Drift Pkts & Acc (\%) & Lat (ms) & Energy (J) \\ \midrule
0.5 & 737k & 3.6k & 93.65 $\pm$ 0.16 & 2.15 $\pm$ 1.26 & 2434 $\pm$ 440 \\
1.0 & 737k & 7.2k & 93.53 $\pm$ 0.19 & 0.99 $\pm$ 0.04 & 2385 $\pm$ 47 \\
5.0 & 737k & 36.9k & 93.59 $\pm$ 0.05 & 0.53 $\pm$ 0.00 & 3675 $\pm$ 47 \\
10.0 & 737k & 74.0k & 93.46 $\pm$ 0.04 & 0.49 $\pm$ 0.01 & 5409 $\pm$ 52 \\
15.0 & 737k & 110.5k & 93.52 $\pm$ 0.05 & 0.46 $\pm$ 0.00 & 6982 $\pm$ 41 \\
50.0 & 737k & 368.6k & 93.52 $\pm$ 0.01 & 0.44 $\pm$ 0.00 & 18598 $\pm$ 33 \\
100.0 & 737k & 737k & 93.58 $\pm$ 0.02 & 0.44 $\pm$ 0.01 & 35352 $\pm$ 352 \\ \bottomrule
\end{tabular}
\end{table}

In Table~\ref{tab:gpu_temp}, GPU temperature remained stable as chaos rate increases.

\begin{table}[H]
\centering
\small
\caption{MI250X Junction Temperature Under Load}
\label{tab:gpu_temp}
\begin{tabular}{cccc}
\toprule
Rate & Runs & Avg (\textdegree C) & Peak (\textdegree C) \\ \midrule
10 & 5 & 43.29 $\pm$ 0.03 & 44.00 $\pm$ 0.00 \\
15 & 5 & 43.29 $\pm$ 0.09 & 44.00 $\pm$ 0.00 \\
50 & 5 & 43.25 $\pm$ 0.05 & 44.00 $\pm$ 0.00 \\
100 & 5 & 43.26 $\pm$ 0.06 & 44.00 $\pm$ 0.00 \\ \bottomrule
\end{tabular}
\end{table}

In Table~\ref{tab:gpu_scalability}, scaling Qwen2.5-7B chaos generation up to 8 MI250X GPUs does not significantly improve latency or energy expenditure.

\begin{table}[H]
\centering
\small
\caption{Qwen2.5-7B Multi-GPU Scalability (MI250X)}
\label{tab:gpu_scalability}
\begin{tabular}{cccc}
\toprule
GPUs & Lat (ms) & Speedup & Energy (J) \\ \midrule
1 & 1973.6 & 1.00 & 12225.7 \\
2 & 2012.5 & 0.98 & 2134.8 \\
3 & 1995.0 & 0.99 & 2117.1 \\
4 & 1988.8 & 0.99 & 2108.0 \\
5 & 2004.5 & 0.98 & 2119.7 \\
6 & 1980.9 & 1.00 & 2105.7 \\
7 & 1991.9 & 0.99 & 2110.2 \\
8 & 1990.7 & 0.99 & 2109.5 \\ \bottomrule
\end{tabular}
\end{table}

\subsection{Multi-GPU Scalability and VRAM Allocation Dynamics}
As documented in Table~\ref{tab:gpu_scalability}, increasing the allocated GPU count from 1 to 8 AMD Instinct MI250X GCDs does not yield a corresponding reduction in inference latency or active energy expenditure. This flat scaling is governed by model footprint and device memory constraints. The Qwen2.5-7B-Instruct model in full bfloat16 precision requires approximately 14.5~GB of memory. Because each individual MI250X GCD is equipped with 64~GB of high-bandwidth memory (HBM2e), the model fits entirely within the memory envelope of a single device. Under default PyTorch/HuggingFace orchestrations without explicit tensor-parallel configurations, the framework prioritizes single-device allocation. Consequently, 100\% of the model parameters are mapped to a single active GPU, leaving the remaining GPUs unutilized. The initial energy spike of 12,225.7~J on 1 GPU is due to the 100.74~s loading and compilation overhead, whereas subsequent multi-GPU runs draw a flat $\sim$2,100~J of baseline energy.

In Table~\ref{tab:gemma_vs_nemotron}, Google's Google Gemma 4 E4B Instruct latency was considerably lower than NVIDIA's Nemotron-4-Mini-15B-Instruct, showing that a smaller, fine-tuned model does experience a significantly lower latency.

\begin{table}[H]
\centering
\small
\caption{Gemma vs. Nemotron Comparative Metrics}
\label{tab:gemma_vs_nemotron}
\begin{tabular}{cccc}
\toprule
Pkt & Gemma Lat (ms) & Nemo Lat (ms) & VQC Selected Reconciler \\ \midrule
W1 & 242.47 & 137.72 & bert \\
W2 & 0.01 & 0.01 & bert \\
W3 & 0.01 & 0.01 & bert \\
1 & 33463.37 & 27422.93 & bert \\
3 & 6302.27 & 6154.50 & bert \\
4 & 1608.36 & 6158.06 & bert \\
7 & 686.74 & 3205.48 & bert \\
8 & 685.48 & 3101.51 & bert \\
9 & 687.45 & 3098.82 & bert \\ \bottomrule
\end{tabular}
\\ \small{\textit{Note: Non-consecutive packet indices (e.g., 2, 5, 6) represent clean telemetry windows that did not experience schema drift and thus bypassed the comparative reconciliation stage. Nemotron results default to Qwen-2.5-7B-Instruct fallback to support offline sandboxed compute execution.}}
\end{table}

\begin{table}[H]
\centering
\small
\caption{VQC Simulator Resilience Metrics ($n=16$, openf1 Stream)}
\label{tab:vqc_sim_resilience}
\begin{tabular}{cccc}
\toprule
API & Acc & Hosseini Index ($C_{\text{rec}}$) & Thru (pps) \\ \midrule
openf1 & 0.934 & 0.934 & 12 \\
openf1 & 0.981 & 0.981 & 11 \\
openf1 & 0.986 & 0.986 & 1 \\ \bottomrule
\end{tabular}
\end{table}

In Table~\ref{tab:qpu_hw_resilience}, and the VQC table, we can see the improvement from the IBM Heron r2 to the VQC quantum routing.

\begin{table}[H]
\centering
\small
\caption{IBM QPU Hardware Resilience Metrics ($n=16$, openf1 Stream)}
\label{tab:qpu_hw_resilience}
\begin{tabular}{cccc}
\toprule
API & Acc & Hosseini Index ($C_{\text{rec}}$) & Thru (pps) \\ \midrule
openf1 & 0.949 & 0.949 & 3 \\
openf1 & 0.946 & 0.946 & 2 \\
openf1 & 0.976 & 0.976 & 1 \\ \bottomrule
\end{tabular}
\end{table}

\section{Discussion of Empirical Anomalies and Baseline Reconciliation}
To ensure empirical consistency and prevent apparent contradictions with prior work (such as the NGISE baseline evaluation), we explicitly clarify the architectural and experimental differences that govern the performance metrics.

\subsection{Latency Discrepancies: Hardware Isolation vs. Broad Sweeps}
In the preliminary NGISE evaluation, the semantic reconciliation layer utilizing the \texttt{all-MiniLM-L6-v2} BERT model documented a baseline latency of approximately $\sim$9.9~ms. However, the 35-run aggregated chaos sweep presented in Table~\ref{tab:chaos_sweep} reports a latency range between 0.44~ms and 2.15~ms under active load. This 5$\times$ to 20$\times$ latency reduction is a direct result of hardware isolation and execution-path optimization. While the NGISE benchmarks were executed across a heterogeneous, un-optimized multi-architecture sweep encompassing consumer-grade CPUs and legacy GPU architectures, the TKDE experiments were conducted on dedicated, water-cooled AMD Instinct MI250X nodes on the LUMI-G supercomputer. These nodes leverage ROCm/HIP graph execution compilation of local execution graphs directly onto CDNA compute units, alongside advanced memory caching of token embeddings.

\subsection{Accuracy Paradox: Structured Telemetry vs. Obfuscated Anomaly Sweeps}
The NGISE evaluation reported an overall mixed accuracy of 62.5\% for the BERT semantic layer, whereas the TKDE chaos sweeps (Table~\ref{tab:chaos_sweep}) maintain a stable accuracy baseline between 93.46\% and 93.65\%. This variation is attributed strictly to the anomaly distribution profile. The NGISE dataset introduced highly destructive, cross-domain anomalies, including heavily obfuscated tokens, randomized abbreviations, and structural truncation. In contrast, the TKDE evaluations isolate the self-healing gateway to structured Formula 1 telemetry parameters. The semantic relationships in F1 streams (e.g., mapping \texttt{velocity\_kph} or \texttt{speed\_mph} to \texttt{speed}) are highly defined and fall within the native semantic alignment of the pre-trained model.

\subsection{Bypass Cache Hit Mechanism (Table 4 Warmups)}
A critical comparison between Table~\ref{tab:chaos_sweep} (BERT latency of 0.44--2.15~ms) and Table~\ref{tab:gemma_vs_nemotron} (Packet indices W2 and W3 registering 0.01~ms latency) reveals an apparent physical impossibility: the same BERT reconciler executing on the same hardware cannot drop in processing latency by two orders of magnitude. This behavior represents a \textit{Bypass Cache Hit} mechanism. The pre-warmup framework utilized in the active telemetry pipeline caches successfully resolved schema mappings. When the gateway encounters a repeated schema modification that has already been reconciled in the current stream, it bypasses the neural network inference step entirely. This hybrid lookup-inference loop shields the GPU from redundant transformer forward passes, achieving near-zero processing overhead for highly repetitive telemetry anomalies.

\subsection{Convergence of Accuracy and Hosseini Index}
As shown in Table~\ref{tab:vqc_sim_resilience} and Table~\ref{tab:qpu_hw_resilience}, the system accuracy ($Acc$) and the Hosseini Recovery Coverage index ($C_{\text{rec}}$) converge to identical values. This equivalence is mathematically expected under our evaluation protocol: because every single detected telemetry anomaly is processed and evaluated by the gateway, the population of drifted fields is identical to the population of tested fields. Thus, the fraction of successfully resolved drifted fields ($Acc$) maps one-to-one to the recovery coverage ratio ($C_{\text{rec}}$).

\end{document}
